\chapter{LITERATURE SURVEY}
\label{ChapterLiteratureSurvey}

Volatility forecasting has been a central topic in financial econometrics for over four decades. This chapter surveys the evolution from classical parametric models through deep learning approaches and regime-switching methodologies, providing the foundation for the proposed Regime-Aware Temporal Fusion Transformer framework.

\section{Literature Collection \& Segregation}

The literature for this project was collected from leading academic databases including IEEE Xplore, ScienceDirect, Springer, SSRN, and Google Scholar. The search focused on three primary domains: (1) volatility forecasting models, (2) deep learning for financial time series, and (3) regime-switching approaches. Over 80 papers were initially identified, from which 22 were selected based on relevance, recency, and methodological contribution.

The selected literature was segregated into three categories:
\begin{itemize}
    \item \textbf{Classical Volatility Models:} Foundational econometric approaches including ARCH, GARCH, and realized volatility models.
    \item \textbf{Deep Learning for Volatility:} Neural network architectures applied to volatility prediction, including LSTM, attention-based RNNs, and Transformers.
    \item \textbf{Regime-Switching and Interpretable Architectures:} Hidden Markov Models, regime-aware strategies, and the TFT architecture.
\end{itemize}

\section{Critical Review of Literature}

\subsection{Classical Volatility Models}

The field of volatility modelling began with the seminal work of Engle \cite{engle1982}, who introduced the Autoregressive Conditional Heteroskedasticity (ARCH) model demonstrating that time-varying variance could be modelled as a function of past squared residuals. Bollerslev \cite{bollerslev1986} extended this into the Generalized ARCH (GARCH) framework, which added lagged conditional variance terms and became the most widely used parametric volatility model in practice.

Recognizing that negative returns tend to increase volatility more than positive returns of equal magnitude---the so-called leverage effect---Nelson \cite{nelson1991} proposed the Exponential GARCH (EGARCH) model. This model uses logarithmic conditional variance, allowing asymmetric effects without requiring non-negativity constraints. Similarly, Glosten, Jagannathan, and Runkle \cite{glosten1993} introduced the GJR-GARCH model with an explicit indicator function for negative shocks.

The availability of high-frequency data enabled the development of realized volatility measures as a non-parametric alternative to GARCH. Parkinson \cite{parkinson1980} proposed an efficient estimator using only daily high and low prices, while Garman and Klass \cite{garman1980} incorporated opening and closing prices for improved efficiency. Rogers and Satchell \cite{rogers1991} developed an estimator robust to drift in the underlying process. Barndorff-Nielsen and Shephard \cite{bns2004} introduced bipower variation as a means to separate the continuous and jump components of return variation.

Corsi's Heterogeneous Autoregressive model of Realized Volatility (HAR-RV) \cite{corsi2009} uses daily, weekly, and monthly lagged RV as regressors in a simple linear framework. Despite its simplicity, HAR-RV remains the dominant benchmark in realized volatility forecasting due to its ability to capture long-memory patterns through the heterogeneous lag structure. In our dataset, HAR-RV achieves $R^2 > 0.98$, setting a demanding benchmark.

\subsection{Deep Learning for Volatility Forecasting}

The application of deep learning to financial volatility forecasting has grown rapidly. Bucci \cite{bucci2020} provided an early systematic comparison showing that LSTM networks outperform traditional GARCH models in daily volatility forecasting, particularly during crisis periods. The ability of LSTMs to capture long-range dependencies through their gating mechanism makes them naturally suited for volatility time series.

Kim and Won \cite{kim2022} advanced this line of research by incorporating attention mechanisms into RNN architectures. Their attention-augmented models better focus on informative lookback periods, thereby improving realized-volatility $R^2$. This work demonstrated that not all historical observations are equally relevant for volatility prediction---some time steps carry significantly more predictive information.

Wu et al. \cite{wu2023} experimented with standard encoder-decoder Transformers for multi-step volatility prediction on US stocks, reporting clear improvements over both HAR-RV and LSTM baselines. Their study highlighted the advantages of self-attention for capturing complex temporal dependencies in volatility dynamics.

Hybrid approaches that combine classical econometric intuition with neural network flexibility have shown particular promise. Zhang et al. \cite{zhang2023} proposed a GARCH-informed neural network which embeds the GARCH variance equation into the inductive bias of the deep network, ensuring that the model respects known volatility dynamics while learning additional non-linear patterns. Rahimikia and Poon \cite{rahimikia2023} incorporated news sentiment scores alongside price features into a machine learning pipeline and demonstrated that textual information carries volatility-relevant value beyond what price data alone provides.

Christensen et al. \cite{christensen2023} conducted a comprehensive survey of machine learning approaches to volatility forecasting, concluding that neural networks and gradient-boosted trees consistently outperform linear benchmarks, with the largest gains observed during periods of market stress.

\subsection{Regime-Switching Models and the TFT Architecture}

Hamilton \cite{hamilton1989} pioneered the application of Markov-switching dynamics to economic time series, providing a principled framework for modelling structural breaks. In the volatility forecasting context, regime-switching models allow different parameterizations for calm, normal, and stressed market conditions. Nystrup et al. \cite{nystrup2021} demonstrated that portfolio strategies driven by HMM-identified regimes significantly improve risk-adjusted returns, motivating the direct integration of regime labels into forecasting models.

The Temporal Fusion Transformer, introduced by Lim et al. \cite{lim2021}, represents a significant advance in interpretable time series forecasting. The architecture combines several key components: Variable Selection Networks that learn instance-specific feature weights, Gated Residual Networks for effective non-linear processing, a seq2seq structure with LSTM encoder-decoder, and multi-head interpretable attention. This combination provides both accurate forecasts and transparent explanations.

Olorunnimbe and Viktor \cite{olorunnimbe2023} applied the TFT to equity price prediction, demonstrating that the variable-selection outputs provide portfolio managers with actionable insights about which factors are driving market dynamics. Chen et al. \cite{chen2024} extended TFT applications to cryptocurrency markets, showing that the architecture transfers well despite the very different microstructure of digital asset markets.

Patel and Sharma \cite{patel2024} examined volatility forecasting in emerging markets using a hybrid regime-aware deep learning framework, while Zhao et al. \cite{zhao2024} proposed an attention-based multi-scale volatility predictor that works directly with high-frequency data, demonstrating the benefits of multi-resolution feature extraction.

\subsection{Summary of Literature Review}

Table \ref{tab:litreview} summarizes the key contributions reviewed in this chapter.

\begin{table}[h!]
\caption{Summary of Key Literature Reviewed}
\label{tab:litreview}
\centering
\small
\begin{tabular}{p{3.5cm}p{2cm}p{6cm}}
\hline
\textbf{Author(s)} & \textbf{Year} & \textbf{Key Contribution} \\
\hline
Engle & 1982 & Introduced ARCH for time-varying variance \\
Bollerslev & 1986 & Generalized ARCH to GARCH framework \\
Hamilton & 1989 & Markov-switching for economic time series \\
Corsi & 2009 & HAR-RV model for realized volatility \\
Bucci & 2020 & LSTM outperforms GARCH for daily RV \\
Lim et al. & 2021 & Temporal Fusion Transformer architecture \\
Kim \& Won & 2022 & Attention-augmented RNNs for volatility \\
Wu et al. & 2023 & Transformer surpasses HAR-RV and LSTM \\
Zhang et al. & 2023 & GARCH-informed neural networks \\
Christensen et al. & 2023 & ML consistently beats linear benchmarks \\
\hline
\end{tabular}
\end{table}

\section{Research Gap}

While significant progress has been made in both deep learning volatility forecasting and regime-switching models, several gaps remain:

\begin{enumerate}
    \item No existing study has applied the TFT architecture specifically to intraday realized volatility forecasting on Indian equities.
    \item The contribution of autoregressive persistence versus genuine feature-driven prediction in deep learning volatility models has rarely been systematically quantified through controlled ablation.
    \item Direct integration of HMM regime labels as conditioning features within the TFT framework has not been explored.
    \item Cross-stock variation in feature-driven predictability remains poorly understood, particularly in emerging market contexts.
\end{enumerate}

This project addresses these gaps through the proposed Regime-Aware Temporal Fusion Transformer framework.
